% GENERATED FILE. Edit canonical lessons, then run npm run generate:books.
% canonical-source-hash: fnv1a64:e422c67ab3cee1fc
% canonical-lessons: PT-C412-motorizada, PT-C412-trotinete, PT-C412-selim, PT-C412-pedalar, PT-C412-parquimetro

\chapter{Em duas rodas}
\label{ch:pt-em-duas-rodas}
\hlchaptermodality{412}

\begin{chapteropening}
\textbf{By the end of this chapter:} \emph{I can talk about getting round on two wheels: the moped, the scooter, the saddle, pedalling and the parking meter.}

Complete the last lesson of chapter 412: I can talk about getting round on two wheels: the moped, the scooter, the saddle, pedalling and the parking meter.
\end{chapteropening}

\section[a motorizada]{a motorizada — the moped}

\label{lesson:PT-C412-motorizada}



\begin{quote}
Before the new one: say the Portuguese for the mileage, then the Portuguese for the excess.
\end{quote}

\subsection*{You'll want to know: a motorizada}
\textbf{a motorizada} — \textquotedblleft{}the moped\textquotedblright{}.

\begin{quote}
O meu avô ia para o trabalho de motorizada.
\end{quote}

\begin{quote}
\emph{My grandfather used to go to work by moped.}
\end{quote}

\emph{A mota} is a bigger motorbike. In Brazil: \emph{a mobilete}.

\begin{cousinweb}
\emph{motorizada} is built on \emph{motor} 'engine', a learned borrowing from Latin \emph{mōtor} 'mover', from \emph{movēre} 'to move' (Lewis \& Short).
\end{cousinweb}

\begin{grammarlens}[title={What you've built}]
One more word for telling, explaining and describing at B1.
\end{grammarlens}

\subsection*{Your turn}
\begin{itemize}
\raggedright
  \item \emph{Say it:} \emph{a motorizada}
  \item \emph{Say it:} \emph{a motorizada}, once more
  \item \emph{Say it:} a sentence of your own with \emph{a motorizada}
  \item \emph{Recall:} say the Portuguese for the mileage, then the Portuguese for the excess, then say \emph{a motorizada} again
\end{itemize}

\begin{tcolorbox}[breakable,colback=teal!4,colframe=teal!35!black,title={Before you move on}]
What does a motorizada mean? (\textquotedblleft{}The moped\textquotedblright{}.) Say it once more.
\end{tcolorbox}
\section[a trotinete]{a trotinete — the scooter (kick or electric)}

\label{lesson:PT-C412-trotinete}



\begin{quote}
Before the new one: say the Portuguese for the excess, then the Portuguese for the moped.
\end{quote}

\subsection*{You'll want to know: a trotinete}
\textbf{a trotinete} — \textquotedblleft{}the scooter (kick or electric)\textquotedblright{}.

\begin{quote}
Muitos turistas andam de trotinete no centro.
\end{quote}

\begin{quote}
\emph{Many tourists ride scooters in the centre.}
\end{quote}

In Brazil: \emph{o patinete}.

\begin{cousinweb}
\emph{trotinete} was borrowed from French \emph{trottinette}, built on \emph{trottiner} 'to trot along in small steps' with the small-size ending \emph{-ette} (EtymDB).
\end{cousinweb}

\begin{grammarlens}[title={What you've built}]
One more word for telling, explaining and describing at B1.
\end{grammarlens}

\subsection*{Your turn}
\begin{itemize}
\raggedright
  \item \emph{Say it:} \emph{a trotinete}
  \item \emph{Say it:} \emph{a trotinete}, once more
  \item \emph{Say it:} a sentence of your own with \emph{a trotinete}
  \item \emph{Recall:} say the Portuguese for the excess, then the Portuguese for the moped, then say \emph{a trotinete} again
\end{itemize}

\begin{tcolorbox}[breakable,colback=teal!4,colframe=teal!35!black,title={Before you move on}]
What does a trotinete mean? (\textquotedblleft{}The scooter (kick or electric)\textquotedblright{}.) Say it once more.
\end{tcolorbox}
\section[o selim]{o selim — the saddle (of a bicycle)}

\label{lesson:PT-C412-selim}



\begin{quote}
Before the new one: say the Portuguese for the moped, then the Portuguese for the scooter.
\end{quote}

\subsection*{You'll want to know: o selim}
\textbf{o selim} — \textquotedblleft{}the saddle (of a bicycle)\textquotedblright{}.

\begin{quote}
O selim da bicicleta está muito alto.
\end{quote}

\begin{quote}
\emph{The bicycle saddle is too high.}
\end{quote}

A horse's saddle is \emph{a sela}.

\begin{cousinweb}
\emph{selim} is a small form of \emph{sela} 'saddle', which is inherited from Latin \emph{sella} 'seat' (EtymDB).
\end{cousinweb}

\begin{grammarlens}[title={What you've built}]
One more word for telling, explaining and describing at B1.
\end{grammarlens}

\subsection*{Your turn}
\begin{itemize}
\raggedright
  \item \emph{Say it:} \emph{o selim}
  \item \emph{Say it:} \emph{o selim}, once more
  \item \emph{Say it:} a sentence of your own with \emph{o selim}
  \item \emph{Recall:} say the Portuguese for the moped, then the Portuguese for the scooter, then say \emph{o selim} again
\end{itemize}

\begin{tcolorbox}[breakable,colback=teal!4,colframe=teal!35!black,title={Before you move on}]
What does o selim mean? (\textquotedblleft{}The saddle (of a bicycle)\textquotedblright{}.) Say it once more.
\end{tcolorbox}
\section[pedalar]{pedalar — to pedal, to cycle}

\label{lesson:PT-C412-pedalar}



\begin{quote}
Before the new one: say the Portuguese for the scooter, then the Portuguese for the saddle.
\end{quote}

\subsection*{You'll want to know: pedalar}
\textbf{pedalar} — \textquotedblleft{}to pedal, to cycle\textquotedblright{}.

\begin{quote}
Pedalamos ao longo do rio todos os domingos.
\end{quote}

\begin{quote}
\emph{We cycle along the river every Sunday.}
\end{quote}

Built on \emph{o pedal}.

\begin{cousinweb}
\emph{pedalar} is built on \emph{pedal}, a learned borrowing from Latin \emph{pedālis} 'of the foot' (EtymDB).
\end{cousinweb}

\begin{grammarlens}[title={What you've built}]
One more word for telling, explaining and describing at B1.
\end{grammarlens}

\subsection*{Your turn}
\begin{itemize}
\raggedright
  \item \emph{Say it:} \emph{pedalar}
  \item \emph{Say it:} \emph{pedalar}, once more
  \item \emph{Say it:} a sentence of your own with \emph{pedalar}
  \item \emph{Recall:} say the Portuguese for the scooter, then the Portuguese for the saddle, then say \emph{pedalar} again
\end{itemize}

\begin{tcolorbox}[breakable,colback=teal!4,colframe=teal!35!black,title={Before you move on}]
What does pedalar mean? (\textquotedblleft{}To pedal, to cycle\textquotedblright{}.) Say it once more.
\end{tcolorbox}
\section[o parquímetro]{o parquímetro — the parking meter}

\label{lesson:PT-C412-parquimetro}



\begin{quote}
Before the new one: say the Portuguese for the saddle, then the Portuguese for to pedal.
\end{quote}

\subsection*{You'll want to know: o parquímetro}
\textbf{o parquímetro} — \textquotedblleft{}the parking meter\textquotedblright{}.

\begin{quote}
Pago o estacionamento no parquímetro.
\end{quote}

\begin{quote}
\emph{I pay for the parking at the meter.}
\end{quote}

Many cities now have an app as well.

\begin{cousinweb}
\emph{parquímetro} is built on English \emph{parking} and \emph{-metro} 'measurer', from Greek \emph{métron} 'measure' (EtymDB).
\end{cousinweb}

\begin{grammarlens}[title={What you've built}]
One more word for telling, explaining and describing at B1.
\end{grammarlens}

\subsection*{Your turn}
\begin{itemize}
\raggedright
  \item \emph{Say it:} \emph{o parquímetro}
  \item \emph{Say it:} \emph{o parquímetro}, once more
  \item \emph{Say it:} a sentence of your own with \emph{o parquímetro}
  \item \emph{Recall:} say the Portuguese for the saddle, then the Portuguese for to pedal, then say \emph{o parquímetro} again
\end{itemize}

\begin{tcolorbox}[breakable,colback=teal!4,colframe=teal!35!black,title={Before you move on}]
What does o parquímetro mean? (\textquotedblleft{}The parking meter\textquotedblright{}.) Say it once more.
\end{tcolorbox}
